% Surface-section revision, 23 October 2026.
% Sections 2--4 are byte-for-byte unchanged from the incoming referee revision.
% Original file preserved as psh_exhaustion(1).tex.
% See the separate referee report and unified diff for the audit trail.
\documentclass[11pt]{amsart}

\usepackage[T1]{fontenc}
\usepackage{lmodern}
\usepackage{amssymb,mathtools}
\usepackage{microtype}
\usepackage{enumitem}
\usepackage{booktabs,tabularx,placeins}
\usepackage{xcolor}
\usepackage[colorlinks=true,linkcolor=blue!55!black,citecolor=blue!55!black,urlcolor=blue!55!black]{hyperref}

\newtheorem{theorem}{Theorem}[section]
\newtheorem{proposition}[theorem]{Proposition}
\newtheorem{lemma}[theorem]{Lemma}
\newtheorem{corollary}[theorem]{Corollary}
\theoremstyle{definition}
\newtheorem{definition}[theorem]{Definition}
\theoremstyle{remark}
\newtheorem{remark}[theorem]{Remark}
\numberwithin{equation}{section}

\newcommand{\C}{\mathbb C}
\newcommand{\R}{\mathbb R}
\newcommand{\PP}{\mathbf P}
\newcommand{\OO}{\mathcal O}
\newcommand{\Tot}{\operatorname{Tot}}
\newcommand{\Bl}{\operatorname{Bl}}
\newcommand{\ddc}{dd^c}
\newcommand{\psh}{\mathrm{psh}}
\newcommand{\Levi}{\mathcal L}
\newcommand{\barp}{\bar\partial}
\newcommand{\supp}{\operatorname{supp}}
\newcommand{\base}{\operatorname{pr}}
\newcommand{\Uone}{\mathrm U(1)}

\title[Continuous and smooth plurisubharmonic exhaustions]{Continuous and Smooth
Plurisubharmonic Exhaustions:\\
Threefold Counterexamples and Partial Results for Surfaces}
\author{Zhijian Chen}
\address{Institute of Mathematics, Academy of Mathematics and Systems Science, Chinese Academy of Sciences, Beijing 100190, China}
\email{chenzhijian@amss.ac.cn}
\author{Bo Xiao}
\address{School of Mathematical Sciences, University of Chinese Academy of Sciences, Beijing, China}
\email{xb1994@amss.ac.cn}
\date{October 23, 2026}
\subjclass[2020]{Primary 32F17; Secondary 32Q15, 32L05}
\keywords{Plurisubharmonic exhaustion, K\"ahler manifold, semipositive line bundle,
ruled surface, Hironaka modification, Ueda theory, class VII surface, boundary divisor}

\begin{document}


\begin{abstract}
Starting from Koike's modified Grauert example, we construct a smooth
quasi-projective K\"ahler threefold with a locally Lipschitz
plurisubharmonic exhaustion but no $C^2$ plurisubharmonic exhaustion.
A modification over one point gives a non-K\"ahler threefold with the
same properties. A separate, concise surface section establishes
numerical boundary criteria for nonprojective compactifiers and smooth
rulings, and records further affirmative projective and smooth-boundary
cases. A classification table distinguishes these results from
excluded configurations and cases not treated by the arguments.
All compactification hypotheses and the recent elliptic-neighborhood
input are stated explicitly. The unrestricted question in complex
dimension two is not resolved.
\end{abstract}

\maketitle

\section{Introduction and main results}\label{sec:intro}

The regularity of a plurisubharmonic exhaustion matters even for
K\"ahler manifolds. This paper constructs a smooth quasi-projective
threefold admitting a continuous plurisubharmonic exhaustion but no
$C^2$ one, and then inserts a non-K\"ahler configuration over a single
point without changing these exhaustion properties. In dimension two,
the same line-bundle construction cannot give a counterexample. We
prove this obstruction and several affirmative regularization results
for surfaces, and specify the boundary problem left by our reductions.
Throughout the introduction, ``psh'' abbreviates plurisubharmonic, and
``smooth'' means $C^\infty$.

\subsection{Threefold counterexamples and their metric construction}

Our compact base is the modified Grauert ruled surface appearing in
Koike's work \cite[Example~4.1]{Koike2021}. Its failure of semipositivity
is already known; the first-order obstruction, including its $C^2$
version, is not claimed as a new result here
\cite[Theorem~1.4 and Example~4.1]{Koike2021},
\cite[Remark~2.8]{Wu2026}. We include a direct proof to make the
regularity threshold used below explicit. The analytic construction
then consists of a cusp correction, a global maximum, and a radius
criterion converting any hypothetical $C^2$ exhaustion into a
forbidden metric. Failure of smooth semipositivity by itself would
not provide a continuous exhaustion.

The cusp estimate also reaches the locally Lipschitz endpoint: a
sufficiently large multiple of $q^{1/2}$ dominates the initial
curvature in the scaled tangent variables. The coefficient is
essential at this endpoint; for smaller exponents shrinking the
neighborhood already suffices.

\begin{theorem}\label{thm:main}
Let $C$ be the smooth projective completion of
\[
 y^2=x^5-x,
\]
let $K=K_C$, and let
\begin{equation}\label{eq:extension-intro}
 0\longrightarrow K\longrightarrow V\longrightarrow\OO_C
 \longrightarrow0
\end{equation}
be a nonsplit holomorphic extension. Let
\[
 S=\PP_{\mathrm{lines}}(V),\qquad
 \pi:S\longrightarrow C,\qquad
 Y=\PP(K)\subset S,
\]
and set
\begin{equation}\label{eq:main-line}
 L=\OO_S(Y)\otimes\pi^*K,\qquad M=\Tot(L^{-1}).
\end{equation}
Then:
\begin{enumerate}[label=\textup{(\roman*)},leftmargin=2.5em]
 \item $S$ is a smooth projective surface and $M$ is a smooth
 quasi-projective K\"ahler threefold;
 \item $L$ admits a continuous Hermitian metric with plurisubharmonic
 local weights, but admits no $C^2$ Hermitian metric with semipositive
 curvature;
 \item $M$ admits a continuous plurisubharmonic exhaustion, but admits
 no $C^2$ plurisubharmonic exhaustion.
\end{enumerate}
The metric and the exhaustion in \textup{(ii)--(iii)} can be chosen
locally Lipschitz, and hence locally of class $C^{0,\gamma}$ for every
$0<\gamma\le1$. No assertion about existence of a $C^1$ exhaustion
is made.
\end{theorem}

The passage from metrics to exhaustions uses a criterion valid for a
holomorphic line bundle $A$ over any compact complex manifold $Z$:
\[
\begin{gathered}
 \Tot(A^{-1})\text{ admits a }C^r\text{ psh exhaustion}\\
 \Longleftrightarrow\\[-2pt]
 A\text{ admits a }C^r\text{ semipositive Hermitian metric},
 \qquad r\in\{2,\infty\}.
\end{gathered}
\]
This is Corollary~\ref{cor:criterion}. A semipositive metric gives the
squared dual norm exhaustion. Conversely, averaging an arbitrary
exhaustion under fiberwise circle rotation and taking a sufficiently
high level yields a radius function whose negative logarithm is a psh
metric weight. Consequently, the obstruction in Theorem~\ref{thm:main}
rules out every $C^2$ psh exhaustion, not only a smoothing of the
particular one constructed here.

The non-K\"ahler example has the same dimension and agrees with $M$
away from a compact exceptional fiber.

\begin{theorem}\label{thm:nonkahler-threefold}
For every point $x$ of the threefold $M$ in Theorem~\ref{thm:main}, there
is a smooth non-K\"ahler threefold $N$ and a proper holomorphic
modification $\mu:N\to M$ inducing a biholomorphism
\[
 N\setminus\mu^{-1}(x)\simeq M\setminus\{x\}.
\]
The threefold $N$ admits a locally Lipschitz psh exhaustion but no
$C^2$ psh exhaustion.
\end{theorem}

The construction places Hironaka's opposite-order blowups inside the
exceptional plane of the blowup of $x$. It produces two compact rational
curves whose sum is zero in real homology, precluding a K\"ahler form.
The exhaustion assertions follow from Proposition~\ref{prop:compact-modification}:
a proper surjective holomorphic map that is biholomorphic away from a compact
subset preserves the existence of $C^r$ psh exhaustions, for
$r=0,2,\infty$. The reverse implication uses convex composition to
make an exhaustion constant near the exceptional set before descending
it. Thus the modification is localized over one point; no small
deformation of a fixed complex structure is asserted. Choosing $x$
outside the zero section also preserves a neighborhood of the compact
base. The compact product principle in Proposition~\ref{prop:compact-product}
gives further examples in higher dimensions, including the Hopf-surface
product in Corollary~\ref{cor:nonkahler}.


\subsection{Selected surface results and their scope}

The two-dimensional statements are collected in Section~\ref{sec:surfaces}.
For a line bundle $A$ on a compact Riemann surface, its dual total
space has a continuous psh exhaustion precisely when $\deg A\ge0$,
and then it has a smooth one (Proposition~\ref{cor:curve-total-space}).
This excludes a curve-base version of our threefold construction, not
all possible surface counterexamples. For arbitrary noncompact
surfaces with a nonconstant holomorphic function, regularization is
a consequence of V\^aj\^aitu's holomorphic-convexity theorem
(Proposition~\ref{prop:surface-nonconstant-function}).

The remaining geometric results assume a divisor compactification
$X=S\setminus\supp D$. For nonprojective $S$, and for a smooth
$\PP^1$-bundle with reduced SNC boundary, continuous and smooth psh
exhaustions are equivalent to the absence of a negative definite
connected boundary component (Theorems~\ref{thm:nonprojective}
and~\ref{cor:ruled-all-degrees}). Further affirmative cases include
disconnected projective boundary, suitable moving square-zero
boundary divisors, torsion canonical bundle, and vertical boundaries
of fibrations (Proposition~\ref{prop:surface-projective-cases}).
The torsion-canonical case is a standard application of surface log
abundance. Proposition~\ref{prop:surface-ueda-sufficient} gives
additional sufficient conditions for smooth square-zero boundaries.

Table~\ref{tab:surface-status} separates what these arguments prove
for compact minimal models from what they leave untreated. Passing
to a minimal compactifier is not automatically compatible with the
open complement. No classification of arbitrary noncompact surfaces,
nor a general continuous-to-smooth regularization theorem in dimension
two, is asserted. Sections~\ref{sec:metrics}--\ref{sec:nonkahler}
contain the threefold constructions; Section~\ref{sec:surfaces} is
independent supplementary material.

\subsection{Relation to prior work and scope of the claims}
\label{subsec:provenance}

The known ruled-surface obstruction, Hironaka's opposite-order blowups,
and the quoted classification and Ueda theorems are inputs, not new
constructions of this paper. The holomorphic-convexity consequence,
the affine-line-bundle construction, and the use of surface abundance
are attributed where used. The surface arguments combine these inputs
with the boundary and character calculations; no priority claim is
made for each intermediate criterion or for a general regularization
theorem.

The self-contained threefold argument is the combination of the
explicit bounded-potential metric and the radius obstruction, followed
by a compactly supported modification. Remark~\ref{rem:minimal}
explains exactly which part of the bounded-potential question in
\cite[Section~4]{Zhang2025} this construction addresses. We make no
claim that dimension three is the smallest possible dimension of a
counterexample, nor that Lipschitz regularity is an optimal threshold.
The full dimension-two problem is not settled in this article.

The elliptic-neighborhood step used in the surface section invokes
Theorem~1.6 of the September 2026 preprint \cite{CFSTZ2026}. This input
is independent of both threefold counterexamples. The classification
table refers to the stated hypotheses, not to a purported complete
classification of all continuously weakly complete surfaces.

\subsection{Conventions and terminology}\label{sec:conventions}

We use
\[
 d^c=\frac{i}{4\pi}(\barp-\partial),\qquad
 \ddc=\frac{i}{2\pi}\partial\barp.
\]
For a holomorphic line bundle $A$ and a local holomorphic frame $e$, a
Hermitian metric is written
\[
 \|e\|_h^2=e^{-\varphi}.
\]
Its curvature is $\Theta_h(A)=\ddc\varphi$, representing $c_1(A)$.
If $e_j=g_{ij}e_i$, the metric weights satisfy
\begin{equation}\label{eq:metric-transition}
 \varphi_j=\varphi_i-\log|g_{ij}|^2.
\end{equation}
A positive continuous Hermitian metric is semipositive in the current
sense when its local weights are continuous psh functions. For a $C^2$
metric this agrees with pointwise semipositivity of its curvature.

For a $C^2$ real-valued function $f$, its unnormalized Levi form is
\[
 \Levi_f(\xi)=\sum_{j,k}f_{z_j\bar z_k}\xi_j\overline{\xi_k}.
\]
Thus $f$ is psh exactly when $\Levi_f\ge0$. The fixed positive factor
$1/(2\pi)$ in $\ddc$ does not affect any positivity argument.
Poincar\'e--Lelong takes the form
\begin{equation}\label{eq:PL}
 \ddc\log\|s_D\|_h^2=[D]-\Theta_h(\OO(D)),
\end{equation}
where $s_D$ is the canonical section of an effective divisor $D$.

\begin{definition}
A continuous real-valued function $F$ on a complex manifold $X$ is an
\emph{exhaustion} if every sublevel set $\{F\le c\}$ is compact.
Equivalently, $F(x_j)\to+\infty$ for every sequence escaping every
compact subset of $X$.
\end{definition}

We say that $X$ is \emph{continuously weakly complete} or
\emph{smoothly weakly complete} when it admits a continuous or smooth
psh exhaustion, respectively. Specifying the regularity avoids the
varying conventions for weak completeness in the literature.
Our projectivization $\PP_{\mathrm{lines}}(E)$ parametrizes
one-dimensional subspaces of the fibers of $E$, and its tautological
line bundle is $\OO_{\PP(E)}(-1)$.

\section{The ruled surface and continuous semipositive metrics}\label{sec:metrics}

\subsection{The ruled surface and the metric obstruction}\label{sec:ruled}

\subsubsection{The curve and the extension}

The polynomial $x^5-x$ has five distinct roots. The double covering
$C\to\PP^1$ is therefore branched at these five roots and at infinity.
Riemann--Hurwitz gives
\[
 2g(C)-2=2(-2)+6=2,
\]
so $g(C)=2$ and $\deg K=2$. Serre duality gives
\[
 \operatorname{Ext}^1(\OO_C,K)=H^1(C,K)
 \simeq H^0(C,\OO_C)^*\simeq\C.
\]
Consequently a nonzero extension class exists. All nonzero classes give
isomorphic extensions after allowing scalar automorphisms of the end
terms, so the isomorphism class of $V$ is unambiguous.

One can also specify the extension through a Dolbeault operator. Choose
a smooth positive $(1,1)$-form $\beta$ on $C$ with
$\int_C\beta=1$, and regard it as a $K$-valued $(0,1)$-form. On the
smooth bundle $K\oplus\OO_C$, set
\begin{equation}\label{eq:dolbeault-extension}
 \barp_V(s,f)=(\barp_Ks+f\beta,\barp f).
\end{equation}
This operator is integrable, since there are no $(0,2)$-forms on a
curve, and defines \eqref{eq:extension-intro}. A splitting would give a
holomorphic lift $(s,1)$ and hence $\barp_Ks=-\beta$. Since $s$ is a
smooth $(1,0)$-form, Stokes' theorem would then imply
$\int_C\beta=0$, a contradiction. The choice of the standard
identification of $K$-valued $(0,1)$-forms with $(1,1)$-forms only changes
this argument by a sign.

\subsubsection{The negative section and the line bundle}

By GAGA, $V$ is algebraic, so $S=\PP_{\mathrm{lines}}(V)$ is a smooth
projective surface. The inclusion $K\hookrightarrow V$ gives a section
$Y=\PP(K)$ of $\pi$. Its normal bundle is
\begin{equation}\label{eq:normal}
 N_{Y/S}\simeq\operatorname{Hom}(K,V/K)\simeq K^{-1}.
\end{equation}
In particular,
\begin{equation}\label{eq:restriction}
 \OO_S(Y)|_Y\simeq K^{-1},\qquad
 L|_Y\simeq\OO_Y,\qquad Y^2=-2.
\end{equation}
The trivialization of $L|_Y$ in \eqref{eq:restriction} is the canonical
one induced by $K^{-1}\otimes K$.

For context, this line bundle is nef and big. If $F$ is a fiber of $\pi$,
then
\[
 L\cdot Y=0,\qquad L\cdot F=1,\qquad L^2=2.
\]
For every irreducible curve $D\ne Y$ dominating $C$,
\[
 L\cdot D=Y\cdot D+2\deg(\pi|_D)>0.
\]
Every vertical irreducible curve is a fiber. Hence $L$ is nef, and the
usual surface criterion $L^2>0$ implies bigness. For completeness, the
latter also follows from Riemann--Roch: $h^2(S,L^m)=0$ for $m$ large,
because $(K_S-mL)$ has negative intersection with a fixed ample divisor,
and therefore
\[
 h^0(S,L^m)\ge\chi(S,L^m)=m^2+O(m).
\]
Moreover, $2L-Y=Y+2\pi^*K$ is ample by Nakai--Moishezon: its square is
$6$, its intersections with $Y$ and a fiber are $2$ and $1$, and its
intersection with every other irreducible curve is positive.

\subsubsection{First-order transition functions}

Choose a finite coordinate cover $\{U_i\}$ of $C$, frames $k_i$ of $K$,
and holomorphic lifts $v_i$ of $1$ in \eqref{eq:extension-intro}.
On overlaps,
\begin{equation}\label{eq:vector-transitions}
 k_j=g_{ij}k_i,\qquad v_j=v_i+c_{ij}k_i.
\end{equation}
The cocycle $\{c_{ij}k_i\}$ represents the nonzero extension class.
Near $Y$, write a line in $V$ as
\[
 [k_i+w_iv_i].
\]
Expressing this line in the $j$-th frame yields
\begin{equation}\label{eq:normal-transition}
 w_j=\frac{g_{ij}w_i}{1-c_{ij}w_i}.
\end{equation}
Shrinking the neighborhoods if necessary, all denominators are nonzero.
Let $d_i$ be the frame of $\OO_S(Y)$ defined by $s_Y=w_id_i$. Then
\[
 \frac{d_j}{d_i}=\frac{w_i}{w_j}
 =\frac{1-c_{ij}w_i}{g_{ij}}.
\]
Thus the frames $\ell_i=d_i\otimes\pi^*k_i$ of $L$ satisfy
\begin{equation}\label{eq:L-transition}
 \ell_j=(1-c_{ij}w_i)\ell_i.
\end{equation}
Their restrictions to $Y$ agree, in accordance with the trivialization
in \eqref{eq:restriction}.

\begin{proposition}\label{prop:no-metric}
The line bundle $L$ admits no $C^2$ Hermitian metric with semipositive
curvature.
\end{proposition}

\begin{proof}
Suppose such a metric exists, with $C^2$ psh weights $\varphi_i$ in the
frames $\ell_i$. By \eqref{eq:L-transition}, the functions
$\varphi_i|_Y$ agree on overlaps. They define a subharmonic function on
the compact connected curve $Y$, hence a constant. Rescaling the metric,
we may assume
\begin{equation}\label{eq:flat-value}
 \varphi_i(z,0)=0.
\end{equation}
In particular $\varphi_{i,z\bar z}(z,0)=0$. The Levi matrix at $(z,0)$
is positive semidefinite. A positive semidefinite Hermitian matrix with
a zero diagonal entry has zero entries in the corresponding row and
column. Therefore
\[
 \varphi_{i,w\bar z}(z,0)=0.
\]
The $C^1$ function
\[
 a_i(z)=\varphi_{i,w}(z,0)
\]
satisfies the Cauchy--Riemann equation and is holomorphic.

Equations \eqref{eq:metric-transition} and \eqref{eq:L-transition} give
\[
 \varphi_j(z,w_j)
 =\varphi_i(z,w_i)-\log|1-c_{ij}(z)w_i|^2,
\]
where the base coordinate change is understood. Differentiating with
respect to $w_i$ at zero and using \eqref{eq:normal-transition}, we obtain
\begin{equation}\label{eq:coboundary}
 g_{ij}a_j=a_i+c_{ij}.
\end{equation}
The holomorphic lifts $v_i'=v_i-a_ik_i$ consequently agree on overlaps:
\[
 v_j'=v_i+c_{ij}k_i-a_jg_{ij}k_i=v_i-a_ik_i=v_i'.
\]
They define a global holomorphic lift of $1$, splitting
\eqref{eq:extension-intro}. This is a contradiction.
\end{proof}

This calculation is the specialized first-order obstruction occurring
in \cite[Theorem~1.4 and Example~4.1]{Koike2021}. Related obstruction
classes, including their applicability to $C^2$ metrics, are discussed
in \cite{Wu2026}. Neither reference is needed for the preceding proof.

\subsection{A local continuous metric near a negative curve}\label{sec:cusp}

We isolate the analytic estimate used in the construction. The unitary
flatness assumption means that a smooth metric of zero curvature exists
on the restricted line bundle.

\begin{lemma}[Cusp correction]\label{lem:cusp}
Let $D$ be a compact connected smooth curve in a complex surface $X$. Suppose
$\deg N_{D/X}<0$, and let $A$ be a holomorphic line bundle such that
$A|_D$ is unitary flat. Choose a smooth metric $h_0$ on $A$ restricting
to a flat metric on $D$, and a smooth metric $h_D$ on $\OO_X(D)$ whose
curvature is strictly negative on $TD$. Put
\[
 q=\|s_D\|_{h_D}^2.
\]
For every $0<a\le\tfrac12$, there is a constant $\Lambda>0$ such that
the local weights of $h_0$ plus $\Lambda q^a$ are strictly psh off $D$
on a sufficiently small neighborhood of $D$ and extend as continuous
psh weights across $D$. They are locally of class $C^{0,2a}$; in
particular they are locally Lipschitz when $a=\tfrac12$.
For $0<a<\tfrac12$ one may take $\Lambda=1$ after shrinking the
neighborhood. At $a=\tfrac12$, a sufficiently large $\Lambda$ suffices.
\end{lemma}

\begin{proof}
The chosen metrics exist. Negative degree on the curve gives a negative
curvature metric on $N_{D/X}=\OO_X(D)|_D$. Any metric on a restricted
line bundle extends smoothly to a neighborhood: extend the logarithm
of its ratio to a fixed smooth metric. The same argument extends the
flat metric on $A|_D$.

Use finitely many relatively compact coordinates $(z,w)$ near $D$,
with $D=\{w=0\}$. Write
\[
 q=|w|^2e^{-b(z,w)},\qquad \beta=-b,
\]
where $b$ is the smooth weight of $h_D$. Negativity along $TD$ says
\begin{equation}\label{eq:beta-positive}
 \beta_{z\bar z}(z,0)>0.
\end{equation}
For $w\ne0$, set $\tau=\log q=\log|w|^2+\beta$. The chain rule gives
the exact identity
\begin{equation}\label{eq:cusp-exact}
 \Levi_{q^a}(v)
 =q^a\bigl(a\Levi_{\tau}(v)+a^2|\partial\tau(v)|^2\bigr).
\end{equation}
Evaluate this on
\[
 v=\xi\frac{\partial}{\partial z}
    +w\nu\frac{\partial}{\partial w}.
\]
Since $\log|w|^2$ is pluriharmonic off $D$,
\[
 \Levi_{\tau}(v)
 =\beta_{z\bar z}(z,0)|\xi|^2
   +O(|w|)(|\xi|^2+|\nu|^2),
\]
and
\[
 \partial\tau(v)
 =\nu+\beta_z(z,0)\xi+O(|w|)(|\xi|+|\nu|).
\]
Therefore
\begin{align}\label{eq:scaled-cusp}
 \frac{\Levi_{q^a}(\xi,w\nu)}{q^a}
 &=a\beta_{z\bar z}(z,0)|\xi|^2
   +a^2|\nu+\beta_z(z,0)\xi|^2\notag\\
 &\hspace{1em}+O(|w|)(|\xi|^2+|\nu|^2).
\end{align}
All remainder bounds are uniform after choosing the finite cover and
shrinking its coordinate neighborhoods.

The leading Hermitian form in \eqref{eq:scaled-cusp} is positive
definite: vanishing of its first term forces $\xi=0$ by
\eqref{eq:beta-positive}, and vanishing of the second then forces
$\nu=0$. Its smallest eigenvalue has a positive uniform lower bound on
the compact coordinate pieces covering $D$. Also $q^a$ is uniformly
comparable to $|w|^{2a}$. Thus, after shrinking the neighborhoods,
\begin{equation}\label{eq:cusp-lower}
 \Levi_{q^a}(\xi,w\nu)
 \ge c|w|^{2a}(|\xi|^2+|\nu|^2)
\end{equation}
for a uniform $c>0$.

Let $\varphi_0$ be a local weight of $h_0$. Its restriction to $D$ has
zero curvature, so $\varphi_{0,z\bar z}(z,0)=0$. Smoothness gives
\[
 \varphi_{0,z\bar z}=O(|w|),\qquad
 \varphi_{0,z\bar w}=O(1),\qquad
 \varphi_{0,w\bar w}=O(1).
\]
Consequently, in the same scaled variables,
\begin{equation}\label{eq:baseline-bound}
 |\Levi_{\varphi_0}(\xi,w\nu)|
 \le C|w|(|\xi|^2+|\nu|^2).
\end{equation}
Combining \eqref{eq:cusp-lower} and \eqref{eq:baseline-bound} yields
\begin{equation}\label{eq:cusp-coefficient}
 \Levi_{\varphi_0+\Lambda q^a}(\xi,w\nu)
 \ge\bigl(\Lambda c|w|^{2a}-C|w|\bigr)
       (|\xi|^2+|\nu|^2).
\end{equation}
If $0<a<\tfrac12$, take $\Lambda=1$ and then shrink until the
coefficient is positive. If $a=\tfrac12$, fix $\Lambda>C/c$; the
coefficient becomes $(\Lambda c-C)|w|>0$ throughout the punctured
neighborhood where the two uniform estimates hold. The constants
$c,C$ were obtained before choosing $\Lambda$, so this choice is not
circular. Every tangent vector at $w\ne0$ can be written in these
scaled variables. This proves strict plurisubharmonicity off $D$ in
both cases.

The function $q^a=|w|^{2a}e^{-ab}$ is continuous and locally
$C^{0,2a}$, including $2a=1$. The local weights
$\varphi_0+\Lambda q^a$ are bounded near $D$ and extend continuously
to it. The removable singularity theorem for psh functions across an
analytic hypersurface therefore proves that these extended weights
are psh. Adding the globally defined function $\Lambda q^a$ preserves
the metric transition law \eqref{eq:metric-transition}.
\end{proof}

\begin{remark}
The scale in \eqref{eq:scaled-cusp} is essential. The normal component
of the Levi form has order $|w|^{2a-2}$, while the tangential component
has order $|w|^{2a}$. These two scales allow the cusp to dominate the
bounded mixed curvature of a smooth initial metric. At $a=\tfrac12$
the scaled cusp and the initial error both have order $|w|$, which
explains the coefficient choice in \eqref{eq:cusp-coefficient}.
Ordinary smoothing of the cusp would not preserve this domination at
$D$. The argument proves no $C^1$ regularity.
\end{remark}

\subsection{An explicit global continuous metric}\label{sec:global}

Return to the ruled surface $S$ and the line bundle $L$ in
\eqref{eq:main-line}. Choose a positive smooth metric $h_K$ on $K$.
Using \eqref{eq:restriction}, extend the dual metric $h_K^{-1}$ on
$\OO_S(Y)|_Y$ to a smooth metric $h_Y$ on $\OO_S(Y)$. Set
\begin{equation}\label{eq:baseline-metric}
 h_0=h_Y\otimes\pi^*h_K,\qquad
 \theta_0=\Theta_{h_0}(L),\qquad
 q=\|s_Y\|_{h_Y}^2.
\end{equation}
The metric $h_0|_Y$ is the flat metric under the canonical
trivialization of $L|_Y$, and
\begin{equation}\label{eq:curvature-restrictions}
 \theta_0|_{TY}=0,\qquad
 \Theta_{h_Y}(\OO(Y))|_{TY}=-\Theta_{h_K}(K)<0.
\end{equation}

Fix $0<a\le\tfrac12$ and choose $\Lambda>0$ as in
Lemma~\ref{lem:cusp}. Since $q$ vanishes exactly on $Y$ and $S$ is
compact, there exists $\varepsilon>0$ such that the weights
$\varphi_0+\Lambda q^a$ are psh on $\{q<\varepsilon\}$.
Write $Q=\max_Sq$ and choose $0<\varepsilon<Q$. Choose a constant
\begin{equation}\label{eq:B-choice}
 B>\max_{\varepsilon\le t\le Q}\bigl(\Lambda t^a-\log t\bigr).
\end{equation}
Define, on $[0,Q]$,
\begin{equation}\label{eq:chi}
 \chi(0)=0,\qquad
 \chi(t)=\max\{\Lambda t^a,\log t+B\}\quad(t>0),
\end{equation}
and put
\begin{equation}\label{eq:global-metric}
 h=h_0\exp\bigl(-\chi(q)\bigr).
\end{equation}

\begin{proposition}\label{prop:continuous-metric}
The metric \eqref{eq:global-metric} is a positive continuous Hermitian
metric on $L$ with psh local weights. It is locally $C^{0,2a}$,
and is locally Lipschitz for $a=\tfrac12$.
\end{proposition}

\begin{proof}
The logarithmic branch in \eqref{eq:chi} tends to $-\infty$ as $t\to0$,
so $\chi(t)=\Lambda t^a$ for all sufficiently small $t$. Hence $\chi$ is
continuous at zero and bounded on $[0,Q]$.

On $S\setminus Y$, Poincar\'e--Lelong and
\eqref{eq:baseline-metric} give
\begin{equation}\label{eq:log-curvature}
 \theta_0+\ddc\log q=\pi^*\Theta_{h_K}(K)\ge0.
\end{equation}
As a current on $S$, the same expression equals
$\pi^*\Theta_{h_K}(K)+[Y]$. Thus $\varphi_0+\log q+B$ is a singular
psh weight of $L$.

On $\{0<q<\varepsilon\}$, the weights of $h$ are
\[
 \varphi_0+\chi(q)
 =\max\{\varphi_0+\Lambda q^a,\varphi_0+\log q+B\},
\]
a maximum of two psh functions. They are therefore psh. On
$\{q\ge\varepsilon\}$, \eqref{eq:B-choice} makes the logarithmic
branch strictly dominant. This domination also holds on a collar of
$\{q=\varepsilon\}$, so no boundary gluing issue arises. The weights
are psh there by \eqref{eq:log-curvature}. Near $Y$, the cusp branch is
dominant, and Lemma~\ref{lem:cusp} proves positivity across $Y$.

All the weights are finite and continuous. Near $Y$ their only
nonsmooth term is $\Lambda q^a=\Lambda|w|^{2a}e^{-ab}$,
which is $C^{0,2a}$ for $0<2a\le1$.
Away from $Y$, both branches are smooth, and their maximum is locally
Lipschitz, hence locally $C^{0,2a}$. This proves the asserted regularity.
\end{proof}

\begin{remark}[Bounded potentials and minimal singularities]\label{rem:minimal}
The current
\[
 T=\theta_0+\ddc\chi(q)\in c_1(L)
\]
has locally Lipschitz potentials when $a=\tfrac12$. It also has minimal
singularities: if $T'=\theta_0+\ddc v\ge0$ is another positive current
in $c_1(L)$, its quasi-psh potential $v$ is bounded above on compact
$S$, while $\chi(q)$ is bounded below. Hence
\[
 v\le\chi(q)+C_v
\]
for a constant $C_v$. This is precisely the comparison expressing that
$T$ is no more singular than $T'$. Here global potentials exist by the
$\partial\bar\partial$-lemma on the projective surface.

For this modified Grauert example, Proposition~\ref{prop:continuous-metric}
therefore gives an affirmative answer to the bounded-potential question
in \cite[Section~4]{Zhang2025} for the modified example. That question
also mentions Grauert's original example; the present argument does
not identify its line bundle with $L$ and does not settle that other
case. The metric is constructed here, not taken from that preprint.
This explicit answer for $L$ is not a statement about arbitrary nef
and big line bundles, nor a claim that every minimal-singularity
representative has continuous potentials.
\end{remark}

\section{Line-bundle exhaustions and the K\"ahler counterexample}\label{sec:radius}

\subsection{Dual norms and compact-group averaging}

\begin{lemma}\label{lem:norm}
Let $Z$ be a compact complex manifold and $A\to Z$ a holomorphic line
bundle with a positive continuous Hermitian metric $h$ having psh
local weights. Then the squared dual norm
\[
 \rho(v)=\|v\|_{h^{-1}}^2\qquad(v\in\Tot(A^{-1}))
\]
is a continuous psh exhaustion. If the weights are locally
$C^{0,\gamma}$ for $0<\gamma\le1$, so is $\rho$.
\end{lemma}

\begin{proof}
For a local frame $e_i$ of $A$, write a dual vector as
$v=\lambda_i e_i^*$, and write $h(e_i,e_i)=e^{-\varphi_i}$. Then
\begin{equation}\label{eq:local-norm}
 \rho=|\lambda_i|^2e^{\varphi_i}.
\end{equation}
Off the zero section, $\log\rho=\log|\lambda_i|^2+\varphi_i$ is psh.
Composition with the increasing convex function $\exp$ makes $\rho$
psh there. Formula \eqref{eq:local-norm} extends continuously by zero
to the zero section and is locally bounded, so removability gives psh
on the whole total space.

Fix a smooth Hermitian metric $h_{\mathrm{ref}}$ on $A^{-1}$. By
compactness of $Z$ and positivity of $h$, constants $c,C>0$ satisfy
\[
 c\|v\|_{h_{\mathrm{ref}}}^2
 \le\rho(v)\le C\|v\|_{h_{\mathrm{ref}}}^2.
\]
The closed disk bundle of any fixed radius over $Z$ is compact.
Every sublevel of $\rho$ is a closed subset of such a disk bundle and
is therefore compact. The H\"older assertion follows locally from
\eqref{eq:local-norm}.
\end{proof}

\begin{lemma}[Averaging preserves exhaustions]\label{lem:average}
Let a compact group $G$ act continuously on a locally compact Hausdorff
space $X$, and let $F$ be a continuous exhaustion. Its normalized Haar average
\[
 \overline F(x)=\int_G F(gx)\,dg
\]
is also an exhaustion. If $G$ is a compact Lie group, $X$ is a complex
manifold, the action is smooth and holomorphic, and $F$ is $C^2$ psh,
then $\overline F$ is $C^2$ psh. The same assertion holds in the
smooth category.
\end{lemma}

\begin{proof}
For $K_c=\{F\le c\}$, the orbit saturation $G\cdot K_c$ is compact.
If $x\notin G\cdot K_c$, then $F(gx)>c$ for every $g\in G$, and
continuity on the compact orbit implies $\overline F(x)>c$. Hence
\[
 \{\overline F\le c\}\subset G\cdot K_c.
\]
The sublevel on the left is closed and therefore compact. Differentiation
under the integral preserves $C^2$ regularity, and averaging the
nonnegative Levi forms of the holomorphic pullbacks preserves psh.
\end{proof}

\subsection{The radius criterion}

\begin{proposition}[The radius criterion]\label{prop:radius}
Let $Z$ be a compact complex manifold and $A\to Z$ a holomorphic line
bundle. If $\Tot(A^{-1})$ admits a $C^2$ psh exhaustion, then $A$ admits
a $C^2$ Hermitian metric with semipositive curvature. If the exhaustion
is smooth, the metric is smooth.
\end{proposition}

\begin{proof}
Average the exhaustion over the holomorphic circle action
$v\mapsto e^{i\alpha}v$ on the fibers, using
Lemma~\ref{lem:average}. Denote the resulting invariant exhaustion by
$F$. For a local frame $e_i$ of $A$ and $v=\lambda_i e_i^*$, circle
invariance gives
\begin{equation}\label{eq:radial-function}
 F(z,\lambda_i)=U_i(z,t),\qquad t=\log|\lambda_i|^2,
 \quad\lambda_i\ne0.
\end{equation}
The function
\begin{equation}\label{eq:holomorphic-pullback}
 V_i(z,\zeta)=U_i(z,\operatorname{Re}\zeta)
             =F(z,e^{\zeta/2})
\end{equation}
is a holomorphic pullback of $F$ and is psh. In particular
$U_{i,tt}\ge0$.

For fixed $z$, continuity at the zero section gives the finite limit
\[
 \lim_{t\to-\infty}U_i(z,t)=F(z,0).
\]
Exhaustiveness gives $U_i(z,t)\to+\infty$ as $t\to+\infty$.
A differentiable convex function with finite left limit has
nonnegative derivative. Moreover, if its derivative vanishes at
$t_0$, convexity forces it to be constant on $(-\infty,t_0]$.
Consequently
\begin{equation}\label{eq:radial-positive}
 U_{i,t}(z,t)>0\quad\text{whenever }U_i(z,t)>F(z,0).
\end{equation}

Fix $c>\max_ZF|_Z$. By convexity, the two limits, and
\eqref{eq:radial-positive}, there is a unique solution
\begin{equation}\label{eq:implicit-radius}
 U_i(z,\sigma_i(z))=c,
 \qquad U_{i,t}(z,\sigma_i(z))>0.
\end{equation}
The implicit function theorem makes $\sigma_i$ of class $C^2$, or
smooth when $F$ is smooth.

We verify the sign of its Levi form explicitly. Differentiating
\eqref{eq:implicit-radius} twice in base coordinates gives
\begin{align}\label{eq:implicit-hessian}
 0={}&U_{i,z_j\bar z_k}
 +U_{i,z_jt}\sigma_{i,\bar z_k}
 +U_{i,t\bar z_k}\sigma_{i,z_j}\notag\\
 &+U_{i,tt}\sigma_{i,z_j}\sigma_{i,\bar z_k}
 +U_{i,t}\sigma_{i,z_j\bar z_k},
\end{align}
with all derivatives of $U_i$ evaluated at $(z,\sigma_i(z))$.
For $V_i$ in \eqref{eq:holomorphic-pullback}, its $\zeta\bar\zeta$
entry is $U_{i,tt}/4$, and its mixed entries have a factor $1/2$.
Thus \eqref{eq:implicit-hessian} is precisely the identity
\begin{equation}\label{eq:radius-Levi}
 \Levi_{V_i}\bigl(\xi,2\partial\sigma_i(\xi)\bigr)
 =-U_{i,t}\Levi_{\sigma_i}(\xi).
\end{equation}
The left side is nonnegative and $U_{i,t}>0$, so
$\psi_i=-\sigma_i$ is psh.

If $e_j=g_{ij}e_i$, the dual coordinates satisfy
$\lambda_j=g_{ij}\lambda_i$. The level in
\eqref{eq:implicit-radius} consequently transforms by
\[
 \sigma_j=\sigma_i+\log|g_{ij}|^2,
\]
so
\[
 \psi_j=\psi_i-\log|g_{ij}|^2.
\]
These are exactly the metric transition laws
\eqref{eq:metric-transition}. Hence $\|e_i\|^2=e^{-\psi_i}$ defines
the desired semipositive Hermitian metric on $A$.
\end{proof}

\begin{corollary}\label{cor:criterion}
For a holomorphic line bundle $A$ over a compact complex manifold,
$\Tot(A^{-1})$ admits a $C^2$ psh exhaustion if and only if $A$ admits
a $C^2$ semipositive Hermitian metric. The same equivalence holds with
$C^2$ replaced by $C^\infty$.
\end{corollary}

\begin{proof}
The forward implication is Proposition~\ref{prop:radius}. For the
reverse implication, the squared dual norm in Lemma~\ref{lem:norm}
has the regularity of the given metric, by its local expression
\eqref{eq:local-norm}.
\end{proof}

\subsection{Proof of the K\"ahler counterexample}\label{sec:conclusion}

\begin{proof}[Proof of Theorem~\ref{thm:main}]
The construction in Subsection~\ref{sec:ruled} gives a smooth projective
surface $S$ and an algebraic line bundle $L$. The total space $M$ is
the open subset
\begin{equation}\label{eq:compactification}
 M=\PP_{\mathrm{lines}}(\OO_S\oplus L^{-1})
       \setminus\PP_{\mathrm{lines}}(L^{-1}).
\end{equation}
The identification sends $v\in L^{-1}_z$ to the line $[(1,v)]$.
The projective bundle in \eqref{eq:compactification} is a smooth
projective threefold. Restricting any of its K\"ahler forms proves
that $M$ is K\"ahler.

Let $h$ be the continuous semipositive metric from
Proposition~\ref{prop:continuous-metric}. If $p:M\to S$ is the bundle
projection, the exhaustion is explicitly
\begin{equation}\label{eq:explicit-exhaustion}
 \rho(v)=\|v\|_{h^{-1}}^2
       =\|v\|_{h_0^{-1}}^2
          \exp\!\bigl(\chi(q(p(v)))\bigr).
\end{equation}
Lemma~\ref{lem:norm} shows that $\rho$ is continuous, psh, and
exhaustive. Taking $a=\tfrac12$ and the coefficient $\Lambda$ supplied
by Lemma~\ref{lem:cusp} makes $h$, its weights, and $\rho$ locally
Lipschitz. In particular the same example has local $C^{0,\gamma}$
regularity for all $0<\gamma\le1$.

If a $C^2$ psh exhaustion on $M$ existed,
Proposition~\ref{prop:radius} would produce a $C^2$ semipositive
Hermitian metric on $L$. Proposition~\ref{prop:no-metric} excludes
such a metric. Therefore no $C^2$ psh exhaustion exists, completing
the proof.
\end{proof}

\begin{remark}[Interpretation]
The example separates existence of a continuous exhaustion from
existence of any smooth exhaustion. The conclusion does not depend on
whether the particular function \eqref{eq:explicit-exhaustion} can be
smoothed: every hypothetical $C^2$ exhaustion leads to the forbidden
metric. The argument does not determine whether a $C^1$ psh
exhaustion exists, and the surface results below do not establish that
dimension three is minimal. Thus a definition of weak pseudoconvexity
requiring a smooth
psh exhaustion is strictly stronger than the corresponding definition
requiring only a continuous psh exhaustion, even for smooth
quasi-projective K\"ahler threefolds.
\end{remark}

\section{Compact modifications and non-K\"ahler counterexamples}
\label{sec:nonkahler}

The K\"ahler counterexample can be changed over a single point to obtain
a non-K\"ahler counterexample of the same dimension.  We first isolate
the exhaustion argument, then insert Hironaka's opposite-order blowup
configuration in an exceptional plane.  Compact products give a simpler
alternative when an increase in dimension is allowed.

\subsection{Exhaustions under compactly supported modifications}

\begin{proposition}\label{prop:compact-modification}
Let $f:X'\to X$ be a proper surjective holomorphic map between nonempty
complex manifolds.  Suppose that, for a compact set $K\subset X$, it
restricts to a biholomorphism
\[
 X'\setminus f^{-1}(K)\longrightarrow X\setminus K.
\]
For every $r\in\{0,2,\infty\}$, $X$ admits a $C^r$ psh exhaustion
if and only if $X'$ does.
\end{proposition}

\begin{proof}
The case $K=\varnothing$ is immediate.  Holomorphic pullback preserves
plurisubharmonicity and the stated
regularity, and properness makes the sublevels of a pulled-back
exhaustion compact.  Conversely, suppose $F$ is a $C^r$ psh exhaustion
on $X'$.  Choose a relatively compact open neighborhood $U$ of $K$ and
\[
 c>\max_{f^{-1}(\overline U)}F.
\]
Choose a smooth convex nondecreasing function
$\chi:\R\to[0,\infty)$ that vanishes on $(-\infty,c]$ and tends to
infinity.  For example, integrate a smooth nondecreasing function
which is zero up to $c$ and one from $c+1$ onward.  On $X\setminus K$
define
\[
 G=\chi\circ F\circ(f|_{X'\setminus f^{-1}(K)})^{-1}.
\]
This function is $C^r$ and psh.  It vanishes on $U\setminus K$, so
extension by zero on $K$ gives a $C^r$ psh function on all of $X$.
For each nonempty sublevel $\{G\le a\}$ choose $b$ such that
$\chi(t)\le a$ implies $t\le b$.  Then
\[
 \{G\le a\}\subset K\cup f(\{F\le b\}).
\]
The left side is closed and the right side is compact.  Thus $G$ is an
exhaustion; no smooth descent of the original function $F$ is needed.
\end{proof}

\subsection{A non-K\"ahler modification over one point}

\begin{proposition}\label{prop:point-modification}
Let $X$ be a connected smooth complex threefold and $x\in X$.
There is a connected smooth non-K\"ahler threefold $N$ and a proper
holomorphic modification $\mu:N\to X$ inducing a biholomorphism
\[
 N\setminus\mu^{-1}(x)\simeq X\setminus\{x\}.
\]
\end{proposition}

\begin{proof}
Blow up $x$, obtaining
\[
 b:W=\Bl_xX\longrightarrow X,\qquad E=b^{-1}(x)\simeq\PP^2.
\]
In homogeneous coordinates $[u:v:w]$ on $E$, take
\[
 A=\{w=0\},\qquad B=\{uv=w^2\}.
\]
The line $A$ and smooth conic $B$ meet transversely at exactly
$p=[1:0:0]$ and $q=[0:1:0]$.  Near either point they are the two
coordinate axes in a smooth surface $E$ inside $W$.

Over $W\setminus\{q\}$, first blow up $A\setminus\{q\}$ and then
the strict transform of $B\setminus\{q\}$.  Over
$W\setminus\{p\}$, first blow up $B\setminus\{p\}$ and then the
strict transform of $A\setminus\{p\}$.  All these centers are closed
smooth curves in the respective open manifolds.  On
$W\setminus\{p,q\}$ the centers are disjoint, so the blowups commute
canonically over the base.  Glue the two resulting manifolds through
this isomorphism and write
\[
 \nu:N\longrightarrow W,\qquad \mu=b\circ\nu.
\]
The glued space is Hausdorff: points with distinct images can be
separated in $W$, and points with the same image belong to one of the
Hausdorff open manifolds lying over the two base opens.  It is smooth
of dimension three because every blowup center is smooth.  Properness
is local on the target, so $\nu$ and $\mu$ are proper.  Each blowup
chart is connected and the charts overlap, hence $N$ is connected.
All centers lie in $E$, so $\mu$ is an isomorphism away from $x$.
This is Hironaka's opposite-order construction
\cite[Section~6]{Barlet2015}; see also \cite[Tag~08J1]{StacksHironaka}.

For completeness, the obstruction to a K\"ahler form is entirely
contained in compact exceptional surfaces.  Let $D_A,D_B\subset N$
be the exceptional divisors dominating $A,B$, and let
$\alpha,\beta\in H_2(N,\R)$ be the classes of their general
$\PP^1$ fibers.  The surfaces $D_A,D_B$ are smooth and compact:
locally each is a ruled surface or its blowup at a point, and they
map properly to the compact curves $A,B$.  The fiber divisors of
each map are linearly equivalent, since $A\simeq B\simeq\PP^1$,
and hence are homologous on the corresponding compact surface.

Over $p$, the strict transform of $B$ meets the first exceptional
divisor transversely at one point.  The second blowup therefore
restricts on that divisor to a surface blowup at that point.  Its
fiber over $p$ is the reduced union of the strict transform
$R_p\simeq\PP^1$ of the first exceptional fiber and the new
exceptional fiber $T_p\simeq\PP^1$. Both occur with multiplicity one.
Indeed, take local coordinates $(x,y,z)$ with
$A=\{y=z=0\}$ and $B=\{x=z=0\}$. In the first blowup chart
$z=yt$, the first exceptional divisor is $y=0$ and the second center
is $x=t=0$. Its restriction to $y=0$ is the blowup of the smooth
surface with coordinates $(x,t)$ at the origin. The pullback of the
fiber equation $x=0$ is its strict transform plus the exceptional
curve, each once. The curve $T_p$ is also the fiber of $D_B\to B$
over $p$. Thus
\[
 \alpha=[R_p]+[T_p],\qquad \beta=[T_p].
\]
Over $q$ the order is reversed.  With analogous notation,
\[
 \beta=[R_q]+[T_q],\qquad \alpha=[T_q].
\]
These relations hold in ordinary homology: they arise from fiber
divisors on compact $D_A$ and $D_B$, regardless of compactness of $N$.
Adding them gives
\begin{equation}\label{eq:hironaka-zero-cycle}
 [R_p]+[R_q]=0\quad\text{in }H_2(N,\R).
\end{equation}
If $\omega$ were a K\"ahler form on $N$, closedness and
\eqref{eq:hironaka-zero-cycle} would imply
\[
 0=\int_{R_p}\omega+\int_{R_q}\omega>0,
\]
a contradiction.
\end{proof}

\begin{proof}[Proof of Theorem~\ref{thm:nonkahler-threefold}]
Apply Proposition~\ref{prop:point-modification} to the threefold $M$
of Theorem~\ref{thm:main} at the chosen point $x$.  If $\rho$ is its
continuous psh exhaustion, then
\[
 \rho_N=\rho\circ\mu,\qquad
 \{\rho_N\le a\}=\mu^{-1}(\{\rho\le a\})
\]
is a continuous psh exhaustion of $N$.  A holomorphic map is locally
Lipschitz in smooth coordinate charts, so local Lipschitz regularity
of the exhaustion is preserved.  A $C^2$ psh
exhaustion on $N$ would give one on $M$ by
Proposition~\ref{prop:compact-modification} with $K=\{x\}$, contrary
to Theorem~\ref{thm:main}.
\end{proof}

\begin{remark}[The sense in which the change is local]
The point $x$ can be chosen outside the zero section of $M\to S$.
The construction then preserves a neighborhood of that section and,
by properness, a neighborhood of infinity.  More precisely, for every
neighborhood $U$ of $x$, $\mu$ is an isomorphism over $M\setminus U$.
This is a modification supported over one point; it makes no claim
about a small deformation of the complex structure on a fixed smooth
manifold.
\end{remark}

\subsection{Compact products and an explicit fivefold}\label{sec:products}

\begin{proposition}[Compact product principle]\label{prop:compact-product}
Let $X$ be a nonempty complex manifold and $T$ a nonempty compact
connected complex manifold.  For each $r\in\{0,2,\infty\}$, $X$
admits a $C^r$ psh exhaustion if and only if $X\times T$ does.  If
$T$ is non-K\"ahler, then $X\times T$ is non-K\"ahler.
\end{proposition}

\begin{proof}
Pullback by the projection preserves regularity and
plurisubharmonicity, with compact sublevels
\begin{equation}\label{eq:product-sublevels}
 \{f\circ\operatorname{pr}_X\le a\}=\{f\le a\}\times T.
\end{equation}
Conversely, restriction to the closed holomorphic slice
$X\times\{t_0\}$ preserves regularity, plurisubharmonicity, and
compact sublevels.  A K\"ahler form on the product would restrict to
one on $\{x_0\}\times T$, proving the final assertion.  In fact, the
maximum principle makes every continuous psh function on the product
constant on each compact connected fiber $\{x\}\times T$.
\end{proof}

\begin{corollary}[A non-K\"ahler fivefold]\label{cor:nonkahler}
Let $M$ be the threefold of Theorem~\ref{thm:main} and let
\begin{equation}\label{eq:hopf-surface}
 H=(\C^2\setminus\{0\})/\langle z\mapsto2z\rangle.
\end{equation}
Then $M\times H$ is a smooth non-K\"ahler fivefold admitting a
continuous psh exhaustion but no $C^2$ psh exhaustion.  For every
$0<\gamma<1$ its exhaustion can be chosen locally $C^{0,\gamma}$.
\end{corollary}

\begin{proof}
The dilation action is free and properly discontinuous: if a compact
set lies in $r\le\|z\|\le R$, only finitely many integers $n$ satisfy
$r/R\le2^n\le R/r$.  Polar coordinates identify its quotient smoothly
with
\begin{equation}\label{eq:hopf-topology}
 H\simeq S^3\times\bigl(\R/(\log2)\mathbb Z\bigr)
 \simeq S^3\times S^1.
\end{equation}
Thus $H$ is a compact connected complex surface with
$H^2_{\mathrm{dR}}(H)=0$.  Any K\"ahler form $\omega$ would be exact,
and Stokes' theorem would give $\int_H\omega^2=0$, contradicting
positivity.  Apply Proposition~\ref{prop:compact-product} and
Theorem~\ref{thm:main}.  Explicitly,
\begin{equation}\label{eq:nonkahler-exhaustion}
 \rho_{M\times H}(v,[z])=\rho(v)
\end{equation}
is the required exhaustion, with the same local H\"older regularity.
\end{proof}

\section{Surface results: proved cases and remaining cases}\label{sec:surfaces}

The question in dimension two is whether a continuous psh exhaustion
implies the existence of a $C^2$, or a smooth, psh exhaustion. All
positive results below give a smooth exhaustion. We do not construct a
surface counterexample. We distinguish statements about an arbitrary
noncompact surface $X$ from statements requiring a compactification
$X=S\setminus\supp D$, where $S$ is a smooth compact connected
complex surface and $D\ne0$ is effective.
The Enriques--Kodaira classification concerns the compact surface $S$,
not $X$. In particular, a classification of compact minimal surfaces is
not a regularization theorem for arbitrary open surfaces.
Table~\ref{tab:surface-status} summarizes exactly the cases established
here.

\subsection{Two basic affirmative cases}\label{sec:surface-curves}

\begin{proposition}[A known holomorphic-convexity consequence]
\label{prop:surface-nonconstant-function}
Let $X$ be a connected noncompact complex surface with a continuous
psh exhaustion. If $\OO(X)\ne\C$, then $X$ has a smooth psh exhaustion.
\end{proposition}
\begin{proof}
The holomorphic-convexity theorem of V\^aj\^aitu
\cite[Theorem~1.1]{Vajaitu2022} makes $X$ holomorphically convex.
Its Remmert reduction $r:X\to Y$ is proper with Stein target. A smooth
strongly psh exhaustion of the complex space $Y$ pulls back to a smooth
psh exhaustion of $X$. Here smoothness on $Y$ means local restriction
of smooth ambient functions, so the pullback is smooth even over
singular points of $Y$.
\end{proof}

This is a consequence of the cited theorem, not a new
holomorphic-convexity result. It leaves the case $\OO(X)=\C$ untreated.
Likewise, the regularization theorem of Mongodi--Slodkowski for domains
in a surface already carrying a real-analytic psh exhaustion
\cite[Corollaries~4.4--4.5]{MongodiSlodkowski2020} requires that ambient
hypothesis; one cannot assume it for $X$ itself.

\begin{proposition}[Line-bundle total spaces over a curve]
\label{cor:curve-total-space}
Let $A$ be a holomorphic line bundle on a compact connected Riemann
surface $B$. The following are equivalent:
\[
 \begin{gathered}
 \Tot(A^{-1})\text{ has a continuous psh exhaustion}\\
 \Longleftrightarrow\quad \deg A\ge0\quad\Longleftrightarrow\\
 \Tot(A^{-1})\text{ has a smooth psh exhaustion}.
 \end{gathered}
\]
\end{proposition}
\begin{proof}
Choose a smooth reference metric with curvature $\theta_0$ and an
area form $\omega$ with $\int_B\omega=1$. For $d=\deg A\ge0$, solve
$\ddc f=d\omega-\theta_0$. The resulting metric has curvature
$d\omega$, and its squared dual norm is a smooth psh exhaustion by
Lemma~\ref{lem:norm}. Conversely, the compactification
\[
 \Tot(A^{-1})=\PP_{\mathrm{lines}}(\OO_B\oplus A^{-1})
                 \setminus\PP_{\mathrm{lines}}(A^{-1})
\]
has infinity section with normal bundle $A$. If $\deg A<0$, that
section has negative self-intersection, contradicting the boundary
obstruction in Lemma~\ref{lem:surface-boundary-tools}\textup{(i)} below.
\end{proof}

Thus the line-bundle construction of Theorem~\ref{thm:main} cannot be
lowered to dimension two by replacing its compact base by a curve.
This does not exclude other surface constructions.

\subsection{Boundary tools and permitted modifications}
\label{sec:surface-pair-reduction}

For a divisor complement, a continuous exhaustion tends uniformly to
$+\infty$ at the boundary: each sublevel is a compact subset of $X$,
hence its complement in $S$ is a neighborhood of $\supp D$. We call a
psh function on a punctured boundary neighborhood which diverges in
this way a \emph{divergent collar}.

Blowing up a boundary point preserves $X$ only when the new boundary
is the \emph{full reduced inverse image}, including the exceptional
curve. This permits resolution of the boundary without changing $X$.
Blowups at interior points give proper modifications with compact
exceptional locus, and Proposition~\ref{prop:compact-modification}
preserves both exhaustion properties. By contrast, contracting a curve
which meets but is not contained in the boundary need not induce a
proper map of the complements. We make no unrestricted reduction of
projective pairs to minimal compactifiers.

\begin{lemma}[Boundary obstruction and collar construction]
\label{lem:surface-boundary-tools}
Let $S$ be smooth compact and let $D\ne0$ be effective.
\begin{enumerate}[label=\textup{(\roman*)},leftmargin=2.5em]
\item\label{item:negative-boundary}
If a connected component of $\supp D$ has negative definite
intersection matrix, then $S\setminus\supp D$ has no continuous psh
exhaustion.
\item\label{item:collar-gluing}
If every connected boundary component admits a smooth divergent psh
collar, then $S\setminus\supp D$ has a smooth psh exhaustion.
\item\label{item:matrix-alternatives}
Let $E=\bigcup E_i$ be a connected reduced simple normal crossing
\textup{(SNC)} divisor and let $\lambda$ be the largest eigenvalue of
$I=(E_i\cdot E_j)$. If $\lambda>0$, $E$ admits a smooth strictly psh
divergent collar. If $\lambda=0$, the matrix is negative semidefinite
with one-dimensional kernel, and there are positive integers $a_i$
with
\begin{equation}\label{eq:surface-null-boundary}
 A=\sum_i a_iE_i,\qquad A\cdot E_i=0,\qquad A^2=0.
\end{equation}
When $S$ is projective, this divisor $A$ is nef.
\end{enumerate}
\end{lemma}
\begin{proof}
For \textup{(i)}, Grauert's contraction theorem \cite{Grauert1962}
contracts the indicated component to a point $p$ of a normal surface,
without changing a punctured neighborhood. A continuous psh function
on such a punctured neighborhood is bounded above near $p$. Indeed,
a local finite projection to a bidisc, with $p$ the only point over
$0$, sends it to the fiberwise maximum $v_*$, a continuous psh function
on the punctured bidisc. This maximum is psh off the branch locus;
continuity and local boundedness extend this across the branch locus
away from $0$. On
\[
 \{|z_1|\le r,\ |z_2|=r\}
\]
it has a finite maximum. The maximum principle on each vertical disc
with $z_1\ne0$, followed by continuity at $z_1=0$, gives the same bound
on the punctured smaller bidisc. This contradicts divergence of an
exhaustion at the contracted boundary.

For \textup{(ii)}, choose disjoint smaller boundary neighborhoods $W_E$
compactly contained in the collar neighborhoods and choose $T_E$ larger
than the collar function on $\partial W_E$. A smooth convex
nondecreasing function $\chi_E$, zero on $(-\infty,T_E]$ and tending to
infinity, makes $\chi_E(\Phi_E)$ vanish near $\partial W_E$. Extend by
zero outside $W_E$ and sum over the finitely many boundary components.
The result is smooth and psh, diverges at every boundary point, and is
proper because $S$ is compact.

For \textup{(iii)}, the off-diagonal entries of $I$ are nonnegative and
its graph is connected. Perron--Frobenius applied after a scalar shift
makes $\lambda$ simple with a strictly positive eigenvector. If
$\lambda>0$, rational approximation and clearing denominators give
$A=\sum a_iE_i$ with $a_i>0$ and $A\cdot E_i>0$ for every $i$.
We spell out the metric construction needed here. The restriction
$\OO_S(A)|_E$ is ample. Embed $E$ by a power of this line bundle and
use its Fubini--Study metric. Lift the embedding sections locally to
$S$ and take the appropriate metric root. Their ambient extensions
have positive tangential curvature along the smooth part of $E$;
at a node they have positive curvature in all directions, since the
Zariski tangent space of $E$ there is the ambient tangent space and
the embedding is injective on tangent spaces. Patch these extensions
using a partition of unity relative to a reference metric, keeping
the chosen metrics near the nodes. The restriction to $E$, and hence
the tangential curvature on its smooth part, is unchanged.

Let $q=\|s_E\|^2$ for any smooth metric on $\OO_S(E)$. On the smooth
part of $E$, $\ddc q$ is positive in the normal direction and zero
in the tangent direction; at a node it is zero. Adding $Cq$ to the
patched weights, for sufficiently large $C$, makes the curvature
strictly positive along all of $E$, by the Schur complement and
compactness, and hence in a neighborhood. For this metric $h$,
$-\log\|s_A\|_h^2$ is the required collar.

If $\lambda=0$, simplicity gives a one-dimensional positive kernel,
which has a rational generator because $I$ is integral. Clearing
denominators proves \eqref{eq:surface-null-boundary}. Intersections
with curves outside $E$ are nonnegative by effectiveness; in the
projective case this proves nefness.
\end{proof}

In particular, under the continuous-exhaustion hypothesis a connected
SNC boundary is either already handled by a positive collar, or
supports the weighted square-zero divisor
\eqref{eq:surface-null-boundary}. A negative definite boundary is an
\emph{exclusion}, not a continuous-versus-smooth counterexample.

\subsection{Nonprojective compactifiers}\label{sec:surface-nonprojective}

\begin{theorem}[Nonprojective boundary criterion]\label{thm:nonprojective}
Let $S$ be a smooth compact connected nonprojective surface and let
$D\ne0$ be effective. Then the following are equivalent:
\begin{enumerate}[label=\textup{(\roman*)},leftmargin=2.5em]
\item $S\setminus\supp D$ has a continuous psh exhaustion;
\item $S\setminus\supp D$ has a smooth psh exhaustion;
\item no connected component of $\supp D$ has negative definite
intersection matrix.
\end{enumerate}
No K\"ahler or SNC hypothesis is required. The Hopf branch of the
proof uses the elliptic-neighborhood theorem in the preprint
\cite[Theorem~1.6]{CFSTZ2026}.
\end{theorem}
\begin{proof}
Only \textup{(iii)}$\Rightarrow$\textup{(ii)} needs proof, by
Lemma~\ref{lem:surface-boundary-tools}\textup{(i)}. We separate the
algebraic-dimension cases, using the classification facts in
\cite[Section~2]{FN2005}.

If $a(S)=1$, the algebraic reduction is a holomorphic elliptic
fibration $f:S\to B$ with connected fibers, and every curve is
vertical. For a fiber $F=\sum m_iC_i$, its matrix $I$ satisfies
\begin{equation}\label{eq:surface-fiber-intersection-identity}
 x^{\mathsf T}Ix=-\sum_{i<j}(C_i\cdot C_j)m_im_j
       \left(\frac{x_i}{m_i}-\frac{x_j}{m_j}\right)^2.
\end{equation}
The kernel is the line spanned by $(m_i)$, and every proper principal
submatrix is negative definite. Hence the boundary contains every
fiber support it meets. Thus $X=f^{-1}(B\setminus P)$ for a finite
nonempty set $P$, and a smooth strictly subharmonic exhaustion of the
punctured compact curve pulls back to a smooth psh exhaustion of $X$.
This argument includes reducible and multiple fibers.

Suppose $a(S)=0$. For every line bundle $L$ on a nonprojective surface,
$c_1(L)^2\le0$ \cite[Lemma~2.1]{FN2005}. Every boundary matrix is
therefore negative semidefinite. Hypothesis \textup{(iii)} and the
same Perron--Frobenius argument give an effective divisor $A$ with
positive integral coefficients, full boundary support, and $A^2=0$.
Contract to a smooth minimal model $p:S\to S_0$. At each blowdown
$q:T\to T'$, write $A=q^*A'+bE$, where $A'=q_*A$ and $E^2=-1$.
Then
\[
 0=A^2=(A')^2-b^2,\qquad (A')^2\le0,
\]
so $b=0$ and $A'\ne0$. Consequently $A=p^*A_0$, $A_0^2=0$, and
\[
 \supp A=p^{-1}(\supp A_0).
\]
The restriction $X\to S_0\setminus\supp A_0$ is proper. It is enough
to construct a smooth exhaustion downstairs.

In the K\"ahler case, $S_0$ is a torus or a K3 surface of algebraic
dimension zero. A torus of algebraic dimension zero has no curves:
an irreducible curve would be elliptic (by adjunction and the absence
of rational curves), and its elliptic quotient would give positive
algebraic dimension. On a K3 surface, Riemann--Roch gives
$h^0(\OO(A_0))\ge2$ because $A_0^2=0$ and $-A_0$ is not effective.
Both alternatives contradict $a(S_0)=0$.

In the non-K\"ahler case, $S_0$ is of class~VII. If $b_2(S_0)=n>0$,
Enoki's Main Theorem \cite[p.~454]{Enoki} identifies a nonzero
effective square-zero divisor with a multiple of the distinguished
cycle on an Enoki surface. Its complement is an affine-line bundle
of degree $-n$ over an elliptic curve \cite[(0.4), p.~453]{Enoki}.
Complete the affine fibers to a projective $\PP^1$-bundle. Its infinity
section has square $n>0$, so
Lemma~\ref{lem:surface-boundary-tools}\textup{(iii),(ii)} gives a smooth
psh exhaustion of the affine-line complement. This uses Enoki's
specific square-zero-divisor theorem, not a classification of all
class~VII surfaces.

If $b_2(S_0)=0$, a minimal class~VII surface containing a curve is
elliptic or Hopf \cite[p.~453]{Enoki}; algebraic dimension zero leaves
Hopf. Each irreducible curve has arithmetic genus one, since
$C^2=K_{S_0}\cdot C=0$. It cannot have rational normalization: a map
from $\PP^1$ lifts to $\C^2\setminus\{0\}$ and is constant. Thus the
curves are smooth elliptic and mutually disjoint. For each of them,
finite Ueda type gives a smooth strictly psh divergent collar
\cite[Theorem~2.5]{CLPT2018}. Infinite type gives a unitary flat
neighborhood by \cite[Theorem~1.6]{CFSTZ2026}, and hence the smooth
pluriharmonic collar $-\log|w_i|^2$ in unitary defining functions.
Lemma~\ref{lem:surface-boundary-tools}\textup{(ii)} finishes the proof.
\end{proof}

In particular, for a nonprojective K\"ahler compactifier, a nonempty
boundary complement can be continuously weakly complete only when
$a(S)=1$ and the boundary is a union of whole fiber supports. The
$a(S)=0$ K\"ahler branch is excluded, not an unresolved positive case.

\subsection{Smooth rulings}\label{subsec:smooth-rulings}

\begin{theorem}[Boundary criterion for a smooth ruling]
\label{cor:ruled-all-degrees}
Let $f:S\to B$ be a smooth $\PP^1$-bundle over a compact connected
Riemann surface, and let $D\ne0$ be reduced SNC. Then continuous psh
exhaustion, smooth psh exhaustion, and the absence of a negative
definite connected boundary component are equivalent for $S\setminus D$.
There is no restriction on the degree $D\cdot F$.
\end{theorem}
\begin{proof}
Necessity follows from Lemma~\ref{lem:surface-boundary-tools}. For
sufficiency, positive largest eigenvalue gives a smooth collar. A
holomorphic $\PP^1$-bundle over a compact curve is a projectivized
rank-two bundle; in particular, it is projective and has a section.
For a
connected component $E=\bigcup E_i$ with zero largest eigenvalue, take
$A=\sum a_iE_i$ as in \eqref{eq:surface-null-boundary}. We construct a
smooth exhaustion of $S\setminus E$, whose restriction supplies the
required collar at $E$.

If every component is vertical, $E$ is one whole fiber, since the
fibers of a smooth ruling are irreducible. Pull back an exhaustion
of $B$ minus its image. Otherwise $A\cdot F>0$, so no fiber can be a
component of $E$. Choose a section $C_0$, with $C_0^2=-e$; the
numerical divisor space has basis $C_0,F$. Writing
$A\equiv aC_0+bF$, the equality $A^2=0$ gives $b=ae/2$.
Writing $E_i\equiv d_iC_0+b_iF$, the identities $A\cdot E_i=0$ give
\[
 E_i\equiv d_i\bigl(C_0+(e/2)F\bigr),\qquad
 E_i^2=E_i\cdot E_j=0\quad(i\ne j).
\]
Distinct components are therefore disjoint. Connectedness and the
SNC hypothesis make $E=C$ a single smooth square-zero curve. For
$d=C\cdot F$, adjunction and
$K_S\equiv-2C_0+(2g(B)-2-e)F$ give
\[
 2g(C)-2=d(2g(B)-2).
\]
Riemann--Hurwitz shows that $C\to B$ is finite \'etale.

If $d=1$, the complement is an affine-line bundle. Its linear bundle
is dual to $N_{C/S}$ and therefore has degree zero. We recall the smooth exhaustion construction of
Koike--Ueda \cite[Proposition~2.1]{KoikeUeda2025}. In unitary flat
frames its transitions are $w_i=a_{ij}w_j+b_{ij}(z)$, $|a_{ij}|=1$.
Choose a smooth section and subtract a global section of the linear
bundle so that its $\bar\partial$ is harmonic. On coordinate discs,
choose antiholomorphic primitives $g_i$ of this harmonic form.
Changing the affine coordinates by holomorphic translations then gives
\[
 w_i=a_{ij}w_j+c_{ij},\qquad
 g_i=a_{ij}g_j+c_{ij},
\]
where $c_{ij}$ is constant: the differences of the $g_i$ are both
holomorphic and antiholomorphic. The functions
$\rho_i=|w_i-g_i|^2$ agree and have Levi matrix
\[
 \begin{pmatrix}|(g_i)_{\bar z}|^2&0\\0&1\end{pmatrix}.
\]
They define a smooth psh exhaustion; its sublevels are bounded closed
disc bundles around a smooth section over compact $B$.

For $d\ge2$, a finite \'etale Galois base change splits the boundary
into $d$ disjoint sections. Each has square zero, since the pullback
of $C$ is their disjoint sum and the projection formula applies.
If $d=2$, putting these sections at $0$ and $\infty$ identifies the
complement with $L^\times$ for a degree-zero line bundle. A unitary
flat metric makes $(\log\|v\|^2)^2$ a smooth psh exhaustion.
If $d\ge3$, three labeled sections trivialize the $\PP^1$-bundle;
the others are constant, as maps from compact $B'$ to $\PP^1$
avoiding $0,1,\infty$. The complement is
$B'\times(\PP^1\setminus Q)$ for a finite set $Q$, and a smooth
subharmonic exhaustion of the punctured projective line pulls back.
Average a nonnegative exhaustion over the finite deck group. The
average remains proper, being at least the original exhaustion divided
by the group order, and descends smoothly to $S\setminus C$.
Finally glue the collars at all connected boundary components using
Lemma~\ref{lem:surface-boundary-tools}\textup{(ii)}.
\end{proof}

This theorem uses the actual smooth ruling, not merely birational
ruledness. It does not claim the same result for arbitrary blowups of
a ruled surface or for singular rational fibrations. For $S=\PP^2$,
every nonzero effective divisor is ample, so its complement has a
smooth psh exhaustion directly from a positive metric on $\OO_S(D)$.

\subsection{Further projective cases and the numerical residual class}
\label{subsec:projective-cases}

\begin{proposition}[Projective affirmative cases]
\label{prop:surface-projective-cases}
Let $S$ be smooth projective, let $D\ne0$ be reduced SNC, and assume
$X=S\setminus D$ has a continuous psh exhaustion. Then $X$ has a
smooth psh exhaustion in each of the following cases:
\begin{enumerate}[label=\textup{(\roman*)},leftmargin=2.5em]
\item $D$ is disconnected;
\item on every connected boundary component with zero largest
intersection eigenvalue, a positive integral kernel divisor $A$
satisfies $\kappa(S,A)=1$ or $K_S\cdot A<0$;
\item $K_S$ is torsion in $\operatorname{Pic}(S)$;
\item every component of $D$ is vertical for a morphism
$f:S\to B$ with connected fibers onto a smooth projective curve.
\end{enumerate}
\end{proposition}
\begin{proof}
For \textup{(i)}, write the connected boundary components as
$E_1,\ldots,E_r$, $r\ge2$. The matrix argument provides effective
$A_k$ with support $E_k$ and $A_k^2\ge0$, with strict inequality in
the positive-eigenvalue case. Since $A_k\cdot A_l=0$ for $k\ne l$,
the Hodge index theorem excludes a positive square and makes the
nonzero square-zero nef classes proportional. The proportionality
factors are positive rational numbers, by intersection with an
integral ample class. Choose positive integers $p_k,q_k$ such that
$p_kA_k-q_kA_1$ is numerically trivial, for $k\ge2$. Its real first
Chern class is zero by Hodge index, and the $\partial\bar\partial$-lemma
gives a unitary flat metric $h_k$ on
$\OO_S(p_kA_k-q_kA_1)$. For its canonical meromorphic section
$\sigma_k$, the functions
\[
 \phi_k=\log\|\sigma_k\|_{h_k}^2,\qquad
 \rho=\sum_{k=2}^r\phi_k^2
\]
are smooth on $X$; the $\phi_k$ are pluriharmonic, so $\rho$ is psh.
It tends to infinity at every boundary component, proving properness.

For \textup{(ii)}, a nef divisor of Iitaka dimension one on a surface
is semiample \cite[Lemma~7.3]{Fujino2021}. If $K_S\cdot A<0$,
Riemann--Roch gives
\[
 \chi(\OO_S(nA))=\chi(\OO_S)-\tfrac n2K_S\cdot A\longrightarrow+\infty.
\]
Serre duality and intersection with an ample class give
$H^2(S,\OO_S(nA))=0$ for large $n$. Thus $\kappa(S,A)\ge1$, and
nefness with $A^2=0$ gives $\kappa(S,A)=1$. In either case $A$ is
semiample. A base-point-free multiple defines, after Stein
factorization, a morphism $f:S\to B$ to a smooth curve with connected
fibers and $\OO_S(mA)=f^*L$. The canonical section descends because
$f_*\OO_S=\OO_B$, so $mA=f^*G$ for a nonzero effective divisor $G$
on $B$. Hence
\[
 S\setminus\supp A=f^{-1}(B\setminus\supp G),
\]
with proper restricted map. Pulling back an exhaustion of the
punctured curve gives a collar at $\supp A$. Combine these with the
positive-eigenvalue collars using
Lemma~\ref{lem:surface-boundary-tools}\textup{(ii)}.

For \textup{(iii)}, consider such a square-zero nef $A$. Choose a
small positive rational $\varepsilon$ so that every coefficient of
$\varepsilon A$ is at most one. The log divisor
$K_S+\varepsilon A$ is nef. Log abundance for surfaces
\cite[Theorem~6.1]{Fujino2012} makes it semiample. Killing the torsion
of $K_S$ and clearing denominators shows that $A$ itself is semiample.
Apply the preceding fibration argument. This is a direct consequence
of log abundance on the given pair; no change of the open surface by
a log minimal model program is involved.

For \textup{(iv)}, identity
\eqref{eq:surface-fiber-intersection-identity} holds for the
multiplicity vector of any fiber. Every proper subset of its
components has negative definite intersection matrix. Therefore the
boundary obstruction forces $D$ to be a union of whole fiber supports.
The restricted map $X\to B\setminus P$ is proper for a finite nonempty
set $P$; pull back an exhaustion of the punctured base.
\end{proof}

In \textup{(i)} no SNC assumption is necessary, because boundary
resolution preserves the complement and its number of boundary
components. Part \textup{(iii)} applies, in particular, to the minimal
projective surfaces of Kodaira dimension zero: K3, Enriques, abelian,
and bielliptic surfaces \cite[Chapter~VI]{BHPV2004}. This is a standard
consequence of the cited abundance theorem, not a new classification
result. Part \textup{(iv)} includes the vertical-boundary case of an
elliptic fibration, but does not treat arbitrary horizontal boundaries.

\begin{corollary}[Residual projective pairs]
\label{cor:surface-connected-residual}
If a divisor complement in a smooth projective surface has a continuous
psh exhaustion but no smooth psh exhaustion, then after boundary
resolution it must have connected reduced SNC boundary $D$ supporting
\[
 \begin{gathered}
 A=\sum_i a_iD_i,\quad a_i\in\mathbb Z_{>0},\quad A\text{ nef},
 \quad A\cdot D_i=0,\quad A^2=0,\\
 \kappa(S,A)=0,\qquad K_S\cdot A\ge0.
 \end{gathered}
\]
The continuous-exhaustion hypothesis is retained; these numerical
conditions alone assert neither its existence nor failure of smooth
exhaustion.
\end{corollary}
\begin{proof}
Disconnected boundary and the positive-eigenvalue case are already
settled. In the remaining case the matrix lemma gives $A$. It is
effective, so $\kappa(S,A)\ge0$, and a nef square-zero divisor has
Iitaka dimension at most one. Proposition~\ref{prop:surface-projective-cases}
excludes dimension one and $K_S\cdot A<0$.
\end{proof}

\subsection{A concise smooth-boundary criterion}
\label{subsec:smooth-boundary}

For a smooth square-zero curve $C$, its normal bundle $N$ is unitary
flat. The Ueda type measures the orders to which local defining
functions can be made to have unitary linear transition functions.
Infinite type is formal; a unitary flat neighborhood requires
\emph{convergent} defining functions $w_i=t_{ij}w_j$, $|t_{ij}|=1$.
These two conditions must not be conflated \cite{Ueda1983,CLPT2018}.

\begin{proposition}[Sufficient conditions, not a classification]
\label{prop:surface-ueda-sufficient}
Let $S$ be a compact K\"ahler surface and let $C\subset S$ be smooth,
connected, and square-zero. Write $\rho:H_1(C,\mathbb Z)\to\Uone$
for the unitary character of $N_{C/S}$ and
$K=\ker[H_1(C,\mathbb Z)\to H_1(S,\mathbb Z)]$.
Then $S\setminus C$ has a smooth psh exhaustion if any of the following
holds: finite Ueda type; a unitary flat neighborhood; torsion or
Diophantine normal bundle; $\rho(K)$ finite; $q(S)\ge g(C)$; or
$g(C)=1$. No continuous exhaustion is assumed. Only the last case
uses \cite[Theorem~1.6]{CFSTZ2026}.
\end{proposition}
\begin{proof}
Finite type gives a smooth strictly psh collar
\cite[Theorem~2.5]{CLPT2018}. A flat neighborhood gives the
pluriharmonic collar $-\log|w_i|^2$. In infinite type, torsion or
Diophantine $N$ gives such a neighborhood by Ueda's convergence
theorem \cite[Theorem~3]{Ueda1983}; finite type has already been
handled. Here Diophantine means
$\operatorname{dist}(N^m,\OO_C)\ge c m^{-\tau}$ for some $c,\tau>0$
and all $m\ge1$.

If $\rho(K)$ is finite, a power $\rho^k$ is trivial on $K$, hence
defines a character of the image of $H_1(C,\mathbb Z)$ in
$H_1(S,\mathbb Z)$. It extends to the whole group because $\Uone$
is divisible. In the remaining non-torsion, infinite-type case,
Neeman's extension theorem in \cite[Theorem~4.4]{CLPT2018} gives an
effective divisor $E$ disjoint from $C$ and cohomologous to $mC$.
A unitary flat metric on $\OO_S(mC-E)$ then gives the pluriharmonic
collar $-m^{-1}\log\|\sigma\|^2$ at $C$.

For $q(S)\ge g(C)$, consider
$r:H^1(S,\OO_S)\to H^1(C,\OO_C)$. If $r$ is not injective,
\cite[Proposition~2.7]{CLPT2018} makes a multiple of $C$ a fiber of a
fibration, and the punctured-base exhaustion applies. Otherwise the
dimension inequality makes $r$ an isomorphism. Hodge decomposition
and duality make $H_1(C,\mathbb Q)\to H_1(S,\mathbb Q)$ an
isomorphism. Its integral kernel is torsion and lies in the free
abelian group $H_1(C,\mathbb Z)$, so $K=0$. Apply the character case.
Finally, when $g(C)=1$, infinite type implies a unitary flat
neighborhood by \cite[Theorem~1.6]{CFSTZ2026}. In every case finish
with Lemma~\ref{lem:surface-boundary-tools}\textup{(ii)}.
\end{proof}

The finite-type and flat-neighborhood arguments also apply without a
K\"ahler hypothesis, as used in the Hopf branch above. Neither
semipositivity of $\OO_S(C)$ nor holomorphic convexity of $S\setminus C$
is asserted as a necessary condition: the proposition constructs
collars, not a general equivalence with either of those properties.

\subsection{What the minimal-surface classification does and does not give}
\label{subsec:surface-classification}

In Table~\ref{tab:surface-status}, $S$ is a \emph{smooth compact minimal}
surface (no smooth rational $(-1)$-curve), $D\ne0$ is reduced SNC, and
$X=S\setminus D$ is assumed to have a continuous psh exhaustion.
``Affirmative'' means that the preceding arguments give a smooth one;
``excluded'' means the hypotheses cannot occur; ``partial'' means
that only the stated additional conditions have been handled. The
table uses the Enriques--Kodaira classification
\cite[Chapter~VI]{BHPV2004}, but every exhaustion assertion is justified
by the results above. It is not a classification of noncompact
surfaces or a claim that every partial case is open in the literature.



The nonprojective theorem already includes nonminimal compactifiers,
by its full-pullback argument. The projective rows do \emph{not}
automatically extend to all surfaces birational to the listed minimal
models. In particular, the table does not settle arbitrary rational
surfaces obtained by boundary-changing blowups of $\PP^2$ or a ruling.
Conversely, an open surface need not be presented by a compact minimal
pair at all. The phrase ``only $\kappa=1,2$ remain'' would therefore
be incorrect for the unrestricted surface problem.

For a smooth connected square-zero boundary on a projective surface,
any putative counterexample must avoid every condition in
Proposition~\ref{prop:surface-ueda-sufficient}; in particular it must
have $g(C)\ge2$, $q(S)<g(C)$, infinite Ueda type without a unitary
flat neighborhood, and $\rho(K)$ infinite. These are necessary
conditions, not examples. The remaining analytic task is to decide
whether a continuous divergent psh collar in such a neighborhood
forces the existence of a smooth one, or to construct a pair where it
does not. The numerical residual conditions do not answer that task.

The general continuous-to-$C^2$ implication for complex surfaces is
not resolved by this paper, even under K\"ahlerness or for all smooth
quasi-projective surfaces. The positive statements in this section,
including the explicitly identified recent elliptic input, are
independent of the threefold counterexamples.

\begin{table}[!htbp]
\caption{Divisor complements with a minimal compactifier; a continuous psh exhaustion is assumed.}
\label{tab:surface-status}
\small
\renewcommand{\arraystretch}{1.16}
\begin{tabularx}{\textwidth}{@{}>{\raggedright\arraybackslash}p{0.29\textwidth}>{\raggedright\arraybackslash}X@{}}
\toprule
Compactifier $S$ & Exhaustion conclusion and exact scope\\
\midrule
Projective, $\kappa(S)=-\infty$: $\PP^2$ or a minimal ruled surface
& \textbf{Affirmative.} Ample boundary in $\PP^2$; the smooth-ruling
criterion otherwise (Theorem~\ref{cor:ruled-all-degrees}).\\
\addlinespace
Projective, $\kappa(S)=0$: K3, Enriques, abelian, bielliptic
& \textbf{Affirmative.} The canonical bundle is torsion;
Proposition~\ref{prop:surface-projective-cases}\textup{(iii)} applies.\\
\addlinespace
Projective, $\kappa(S)=1$: properly elliptic
& \textbf{Partial.} Vertical boundary is handled. Horizontal boundary
requires another stated criterion; no blanket conclusion is proved.\\
\addlinespace
Projective, $\kappa(S)=2$: general type
& \textbf{Partial.} Disconnected, positive-collar, moving, or the
specified smooth-boundary cases are handled. The residual connected
square-zero case is not settled.\\
\midrule
Nonprojective, $a(S)=1$
& \textbf{Affirmative.} Continuous exhaustion forces the boundary to be
a union of whole elliptic fiber supports. This includes the relevant
Kodaira and elliptic surfaces.\\
\addlinespace
Nonprojective K\"ahler, $a(S)=0$: torus or K3
& \textbf{Excluded.} No nonempty divisor complement here satisfies
the continuous-exhaustion hypothesis.\\
\addlinespace
Nonprojective, $a(S)=0$, class~VII, $b_2=0$
& \textbf{Affirmative when a boundary exists.} Such a surface is Hopf;
elliptic collars apply, using \cite{CFSTZ2026}. Curve-free surfaces
have no nonempty divisor boundary.\\
\addlinespace
Nonprojective, $a(S)=0$, class~VII, $b_2>0$
& \textbf{Affirmative for every admissible boundary.} The matrix
condition forces the Enoki square-zero-divisor case. All other
boundaries are excluded by the negative definite obstruction.\\
\bottomrule
\end{tabularx}
\end{table}
\FloatBarrier

\begin{thebibliography}{99}
\raggedright

\bibitem{Koike2021}
T.~Koike,
\emph{Linearization of transition functions of a semi-positive line
bundle along a certain submanifold},
Ann. Inst. Fourier \textbf{71} (2021), no.~5, 2237--2271.
\href{https://doi.org/10.5802/aif.3439}{doi:10.5802/aif.3439}.

\bibitem{Wu2026}
X.~Wu,
\emph{Generalized Ueda Obstruction Classes and Applications to
Non-semi-positivity of Line Bundles},
preprint (2026), arXiv:2606.22821v1.
\href{https://arxiv.org/abs/2606.22821v1}{arXiv:2606.22821}.

\bibitem{Zhang2025}
Y.~Zhang,
\emph{On the Non-Semipositivity of a Nef and Big Line Bundle on
Grauert's Example},
preprint (2025), arXiv:2512.23446v1.
\href{https://arxiv.org/abs/2512.23446v1}{arXiv:2512.23446}.

\bibitem{Vajaitu2022}
V.~V\^aj\^aitu,
\emph{Holomorphic convexity of pseudoconvex surfaces},
Ann. Sc. Norm. Super. Pisa Cl. Sci. (5) \textbf{23} (2022),
no.~4, 1659--1678.
\href{https://doi.org/10.2422/2036-2145.202011_018}{doi:10.2422/2036-2145.202011\_018}.
\url{https://journals.sns.it/index.php/annaliscienze/article/view/5700}.

\bibitem{MongodiSlodkowski2020}
S.~Mongodi and Z.~Slodkowski,
\emph{Domains with a continuous exhaustion in weakly complete surfaces},
Math. Z. \textbf{296} (2020), nos.~3--4, 1011--1019.
\href{https://doi.org/10.1007/s00209-020-02466-z}{doi:10.1007/s00209-020-02466-z}.
\url{https://arxiv.org/abs/1904.03497}.

\bibitem{Grauert1962}
H.~Grauert,
\emph{\"Uber Modifikationen und exzeptionelle analytische Mengen},
Math. Ann. \textbf{146} (1962), 331--368.
\href{https://doi.org/10.1007/BF01441136}{doi:10.1007/BF01441136}.

\bibitem{BHPV2004}
W.~P.~Barth, K.~Hulek, C.~A.~M.~Peters and A.~Van de Ven,
\emph{Compact Complex Surfaces}, second enlarged edition,
Ergebnisse der Mathematik und ihrer Grenzgebiete (3), vol.~4,
Springer-Verlag, Berlin, 2004.
\href{https://doi.org/10.1007/978-3-642-57739-0}{doi:10.1007/978-3-642-57739-0}.

\bibitem{FN2005}
Y.~Fujimoto and N.~Nakayama,
\emph{Compact complex surfaces admitting non-trivial surjective endomorphisms},
Tohoku Math. J. (2) \textbf{57} (2005), no.~3, 395--426.
\href{https://doi.org/10.2748/tmj/1128703004}{doi:10.2748/tmj/1128703004}.
\href{https://www.kurims.kyoto-u.ac.jp/preprint/file/RIMS1431.pdf}{Author preprint, Section~2 and Lemma~2.1}.

\bibitem{Enoki}
I.~Enoki, \emph{Surfaces of class $\mathrm{VII}_0$ with curves}, Tohoku Math. J. (2)
\textbf{33} (1981), no.~4, 453--492.
\href{https://doi.org/10.2748/tmj/1178229349}{doi:10.2748/tmj/1178229349}.
\href{https://www.jstage.jst.go.jp/article/tmj1949/33/4/33_4_453/_pdf/-char/en}{Original article, pp.453--454}.

\bibitem{Fujino2012}
O.~Fujino,
\emph{Minimal model theory for log surfaces},
Publ. Res. Inst. Math. Sci. \textbf{48} (2012), no.~2, 339--371.
\href{https://doi.org/10.2977/PRIMS/71}{doi:10.2977/PRIMS/71}.

\bibitem{Fujino2021}
O.~Fujino,
\emph{Minimal model theory for log surfaces in Fujiki's class $\mathcal C$},
Nagoya Math. J. \textbf{244} (2021), 256--282.
\href{https://doi.org/10.1017/nmj.2020.14}{doi:10.1017/nmj.2020.14}.

\bibitem{Ueda1983}
T.~Ueda,
\emph{On the neighborhood of a compact complex curve with topologically
trivial normal bundle},
J. Math. Kyoto Univ. \textbf{22} (1982/83), no.~4, 583--607.
\href{https://doi.org/10.1215/kjm/1250521670}{doi:10.1215/kjm/1250521670}.

\bibitem{CLPT2018}
B.~Claudon, F.~Loray, J.~V.~Pereira and F.~Touzet,
\emph{Compact leaves of codimension one holomorphic foliations on
projective manifolds},
Ann. Sci. \'Ecole Norm. Sup. (4) \textbf{51} (2018), no.~6, 1457--1506.
\href{https://doi.org/10.24033/asens.2379}{doi:10.24033/asens.2379}.
\href{https://arxiv.org/abs/1512.06623}{Author preprint, Theorems~1.3, 2.5, 4.4 and Proposition~2.7}.

\bibitem{CFSTZ2026}
Y.~Chen, S.~Filip, S.~Sun, V.~Tosatti and J.~Zhang,
\emph{Compact K\"ahler surfaces with semipositive anticanonical bundle},
preprint (2026), arXiv:2609.26716v1, submitted September~22, 2026.
\href{https://arxiv.org/abs/2609.26716v1}{arXiv:2609.26716}.

\bibitem{KoikeUeda2025}
T.~Koike and T.~Ueda,
\emph{Plurisubharmonic functions on affine line bundles over compact
K\"ahler manifolds},
Publ. Res. Inst. Math. Sci. \textbf{61} (2025), no.~1, 139--152.
\href{https://doi.org/10.4171/PRIMS/61-1-3}{doi:10.4171/PRIMS/61-1-3}.

\bibitem{Barlet2015}
D.~Barlet,
\emph{Some examples due to H.~Hironaka},
preprint (2015), arXiv:1501.01129v2, Section~6.
\href{https://arxiv.org/abs/1501.01129}{arXiv:1501.01129}.

\bibitem{StacksHironaka}
The Stacks Project Authors,
\emph{The Stacks Project},
Lemma~110.67.1, Tag~08J1, ``Being projective is not local on the base.''
\url{https://stacks.math.columbia.edu/tag/08J1}.

\end{thebibliography}
\end{document}
